% SPDX-License-Identifier: CC-BY-SA-4.0
% Copyright (c) 2026 Luis Alonso
%
% This file is part of luis_alonso/Notas-Japones.
%
% Licensed under Creative Commons Attribution-ShareAlike 4.0 International (CC BY-SA 4.0).
% You may share and adapt this material for any purpose,
% as long as you provide attribution and distribute adaptations under the same license.
%
% License: https://creativecommons.org/licenses/by-sa/4.0/
% Source:  https://git.lalonso.com/luis_alonso/Notas-Japones
%
% Note: Third-party content (if any) is excluded from this license and is marked where used.

\documentclass[a4paper,lualatex,ja=standard,base=10pt]{bxjsarticle}
\usepackage{setup}
\usepackage{ulem}

\newcommand{\rmv}[1]{\textcolor{red}{\sout{#1}}}

\title{\ruby{第|０３|課}{だい||か} LV4}
\author{アロンゾ　ルイス}

\begin{document}
\setcounter{page}{1}

\maketitle{}
\thispagestyle{fancy}
\section{Transformar なります}

Con イ-Adj y se quita la い y se cambia por く.

\begin{itemize}
  \item \ruby{暖}{あたた}か\rmv{い} pasa a \ruby{暖}{あたた}かくなります
  \item 大きい pasa a 大きくなります
\end{itemize}

Con ナ-Adj y sustantivos se cambia な o solo se pone い.

\begin{itemize}
  \item \ruby{元|気}{げん|き}(adjetivo) pasa a \ruby{元|気}{げん|き}になります
  \item しずか (adjetivo) pasa a しずかになります
  \item 秋 (sustantivo) pasa a 秋になります
\end{itemize}

\section{Dar respuesta a porque te gusta o no te gusta algo}

\subsection{Sustantivar adjetivos}

Para sustantivar un ajetivo se usa el caracter \textcolor{blue}{の} como
sufijo del adjetivo.

Ej. \ruby{暑|い}{あつ|} pasa a \ruby{暑|い}{あつ|}\textcolor{blue}{の}

にがてです: No me gusta.

Para decir que me gusta el calor: あついのがすきですから。

Para decir que me gusta el frio: さむいの好きですから。

辛いのが好きですから。

静かなのが好きですから。

かんたんなのが好きですから。

しずかになりました。 Volverse silencioso.

%%%

Sustantivar adjetivos.

あついのがにがてですから。 No me gusta el calor o no soy bueno con el calor.

静かなのが好きですから。 Me gusta que sea tranquilo.

\section{Vocabulario}
\medskip

\begin{itemize}
  \item うき Temporada de lluvia.
  \item つゆ Temporada de lluvia en Japón.
  \item かんき Sequia.
  \item すずしい Templado.
  \item ふります Caer precipitar, cuando cae nieve, lluvia, ceniza volcanica.
  \item ぶどうがり Cosechar uvas.
  \item むしあつい Calor humedo.
  \item いめかり Cosechar arroz.
  \item うぐいす Pajaro especifico
  \item にゅうがくしき Entrada a la escuela.
  \item たうえ Siembra.
  \item じょやのがめ Campanadas de año nuevo.
  \item なべりょうり Cocina con nabe.
  \item あじさい Hortencia.
  \item すいかわり Romper la sandia.
  \item たんぼ Campo de arroz.
  \item せみ Cigarra.
  \item ひまがり Girasol.
  \item こおろぎ Grillo.
  \item もみじ Cambio de color de las ojas en otoño.
  \item へいきん Promedio.
  \item しょうがつ Año nuevo.
  \item スキー Ski.
  \item さむい Frio.
  \item 花見 Ver las flores.
  \item たおうえ Recoger arroz.
  \item 月み Ver la luna.
  \item ちちめんちょう Pavo.
  \item きせつ Estacion.
  \item だんだん Que se convierte poco a poco.
  \item うるさい Ruidoso.
  \item ことり Polluelo.
  \item ごうつう Transporte.

  \item こわれます Descomponerse.
  \item なおります Arreglarse.
  \item そつぎょう Graduacion.
  \item さくし Letra de canción.
\end{itemize}

% \section{Ejecicios}
% \medskip
%
% かつどう P.36
%
% \begin{enumerate}
%   \item 冬, 3 月
%   \item 夏, 9 月
%   \item 春, 7 月
%   \item 秋, 12 月
% \end{enumerate}
%
% かつどう P.38
%
% \begin{enumerate}
%   \item 春, e
%   \item 夏, f
%   \item 冬, h
%   \item 秋, g
% \end{enumerate}
%
% りかい P40
% \begin{enumerate}
%   \item a
%   \item b
%   \item b
%   \item a
%   \item a
% \end{enumerate}
%
% Ejercicio de diapositiva 1
%
% \begin{enumerate}
%   \item ちさかったです。/大きくなりました。
%   \item ふべんでした。/べんりになりました。
%   \item たくさんありました。/すくなくなりました。
%   \item ひくかったです/たかくなりました。
% \end{enumerate}
%
% Ejecicio de diapositiva 2
%
% \begin{enumerate}
%   \item お金はありませんから、新しい本かえりませんでした。
%   \item パソコンはべんりですから、しごとはいいです。
%   \item 日曜日ですから、休みます。
%   \item あした、いそがしいですから、かじがあります。
%   \item 時かんがありませんから、しゅくだいがありますから。
% \end{enumerate}
%
% Ejercicio de diapositiva 3
%
% \begin{enumerate}
%   \item あついのがにがてですから、秋がすきです。
%   \item あまいのが好きですから、ケーキを食べました。
%   \item あたたかいのが好きですから、春がいちばん好きです。
%   \item さむいのが好きでから、冬が大好きです。
%   \item からいのがにがてですから、タイりょうりはあまり食べません。
%   \item すずしいのが好きですから、秋だい好きです。
% \end{enumerate}
%
% りかい P.44
%
% \begin{enumerate}
%   \item メリボルンは今、どんなきせつですか。
%   \item どんなふくがいいですか。
% \end{enumerate}
%
% \begin{enumerate}
%   \item a
%   \item a
%   \item b
% \end{enumerate}
%
% \begin{enumerate}
%   \item Tシャツとジーンズ
%   \item セーターやジャケット
% \end{enumerate}

\end{document}
